\section{System Model}
\label{sec:system}

We consider a grant-free uplink from $K$ single-antenna terminals to a
low Earth orbit (LEO) satellite equipped with a uniform planar array (UPA).
The array has $M=M_yM_z$ elements. An access block contains $Q$ affine
frequency division multiplexing (AFDM) pilot frames, each with $N$ useful
samples. Terminal activities, channel geometry, and path coefficients
remain fixed within the block.
The receiver knows the pilot sequences and the system parameters,
whereas terminal activities, channel realizations, inter-frame phase
increments, and receiver noise power are unknown.

Let $\alpha_k\in\{0,1\}$ denote the activity state of terminal $k$.
In the activity model, independent activations with
probability $p$ are conditioned on at least one transmission:
\begin{equation}
 p_{\rm gen}(\bm\alpha)=
 \frac{\mathbf{1}\{\sum_{k=1}^{K}\alpha_k\geq1\}}{Z}
 \prod_{k=1}^{K}p^{\alpha_k}(1-p)^{1-\alpha_k},
 \qquad Z=1-(1-p)^K.
 \label{eq:activity-generation}
\end{equation}
The resulting activity marginal is $\rho=p/Z$, which is supplied to
the receiver.

\subsection{Terrestrial-Satellite Link Model}

Let $\mathbf{s}$ and $\mathbf{v}_{\rm sat}$ denote the satellite position
and velocity, respectively. The position of terminal $k$ is denoted by
$\mathbf{g}_k$.
With carrier frequency $f_c$ and propagation speed $c_0$, define
\begin{equation}
 \mathbf{u}_k=\frac{\mathbf{g}_k-\mathbf{s}}
 {\norm{\mathbf{g}_k-\mathbf{s}}},
 \qquad
 \nu_k=-\frac{f_c}{c_0}\mathbf{v}_{\rm sat}^{\mathsf T}\mathbf{u}_k.
 \label{eq:geometry}
\end{equation}
We use the complex-envelope convention
$s_{\rm RF}(t)=\Re\{x(t)e^{-j2\pi f_ct}\}$; thus, a path contributes
$x(t-\tau(t))e^{j2\pi f_c\tau(t)}$ up to its complex gain.
Equation~\eqref{eq:geometry} consequently gives the Doppler time factor
$e^{j2\pi\nu_kt}$. Envelope delays and array directions are held fixed
over the block, while this carrier phase variation is retained.

After compensation of the common reference delay $\tau_c$, the residual
delay of path $\ell$ is given by
\begin{equation}
 \tau_{k\ell}=\frac{\norm{\mathbf{g}_k-\mathbf{s}}}{c_0}-\tau_c
             +\delta_\ell\tau_{{\rm spread},k},
 \label{eq:path-delay}
\end{equation}
where $\tau_c$ is the propagation time to the closest ground point in the
beam footprint, and $\delta_\ell\tau_{{\rm spread},k}$ is the excess
multipath delay. We adopt a common spatial response and a common Doppler
shift for the paths of each terminal
\cite[Sec.~II-C]{li2023}\cite[Sec.~II]{huang2026}.

Let $\mu_{y,k}=\mathbf e_y^{\mathsf T}\mathbf u_k$ and
$\mu_{z,k}=\mathbf e_z^{\mathsf T}\mathbf u_k$, where $\mathbf e_y$ and
$\mathbf e_z$ are the array axes. For element position
$\mathbf s-n_yd_y\mathbf e_y-n_zd_z\mathbf e_z$, the array response is
\begin{equation}
 [\mathbf{a}_k]_m=
 \exp\!\left[j2\pi\left(
 \frac{d_y}{\lambda}n_y\mu_{y,k}
 +\frac{d_z}{\lambda}n_z\mu_{z,k}\right)\right],
 \quad m=n_z+M_zn_y,\quad \norm{\mathbf{a}_k}^2=M,
 \label{eq:upa}
\end{equation}
where $\lambda=c_0/f_c$, $0\leq n_y<M_y$, and $0\leq n_z<M_z$.
The unknown array response is included in the channel to be estimated.

\subsection{Transmission Scheme}

Let $T_s=1/(N\Delta f)$ be the sampling interval, where $\Delta f$
denotes the subcarrier spacing. The discrete affine Fourier transform
(DAFT) matrix is \cite[Secs.~II-B and III-A]{bemani2022}
\begin{equation}
 [\mathbf{F}_{\rm A}]_{rt}=\frac{1}{\sqrt N}
 e^{-j2\pi(c_2r^2+rt/N+c_1t^2)},
 \qquad r,t=0,\ldots,N-1.
 \label{eq:daft}
\end{equation}
For pilot vector $\mathbf p_{qk}$, the useful time samples are
$\mathbf s_{qk}=\mathbf F_{\rm A}^{\mathsf H}\mathbf p_{qk}$.
A chirp-periodic prefix (CPP) of length $N_{\rm CPP}$ extends these samples as
\begin{equation}
 s_{qk}[t]=e^{j2\pi c_1(t^2-\langle t\rangle_N^2)}
 s_{qk}[\langle t\rangle_N],\qquad -N_{\rm CPP}\leq t<0,
 \label{eq:cpp}
\end{equation}
where $\langle x\rangle_N=x-N\lfloor x/N\rfloor$.
We choose $N_{\rm CPP}T_s\geq\max_{k,\ell}\tau_{k\ell}$ and
$2Nc_1\in\mathbb N_+$. The frame duration is
$T_{\rm frame}=(N+N_{\rm CPP})T_s$.

To account for continuous delays, we use the AFDM waveform model in
\cite[Sec.~II]{huang2026}. Let $w_\xi=\langle\xi\rangle_N$ and
$\zeta_n(w)=\lfloor2c_1w+n/N\rfloor$. The basis is
\begin{equation}
 \varphi_n(\xi)=\frac{1}{\sqrt N}
 e^{j2\pi[c_2n^2+c_1\xi^2+nw_\xi/N-\zeta_n(w_\xi)w_\xi]}.
 \label{eq:continuous-waveform}
\end{equation}
It agrees with the inverse DAFT at integer useful-sample times and
extends across prefix periods according to \eqref{eq:cpp}.
For a delay $d=\tau/T_s$, set $w=\langle t-d\rangle_N$ and
$\zeta_{tn}=\zeta_n(w)$. The delayed basis can be expressed as
\begin{equation}
 [\mathbf B_d]_{tn}=\varphi_n(t-d)
 =\frac{1}{\sqrt N}
 e^{j2\pi[c_2n^2+c_1(t-d)^2+nw/N-\zeta_{tn}w]}.
 \label{eq:continuous-basis}
\end{equation}
The corresponding DAFT-domain path operator, for useful-window start
index $g$, is given by
\begin{equation}
 \mathbf G(d,\nu;g)=\mathbf F_{\rm A}\mathbf D_{\nu,g}\mathbf B_d,
 \qquad
 \mathbf D_{\nu,g}=\diag\{e^{j2\pi\nu T_s(g+t)}\}_{t=0}^{N-1}.
 \label{eq:path-operator}
\end{equation}
Expanding the path operator gives
\begin{equation}
 \begin{split}
 G_{rn}
 &=\frac{e^{j2\pi[c_1d^2+c_2(n^2-r^2)-nd/N+\nu T_sg]}}{N}\\
 &\quad\times\sum_{t=0}^{N-1}
 e^{j2\pi[(n-r)/N-2c_1d+\nu T_s]t}
 e^{j2\pi\zeta_{tn}d}.
 \end{split}
 \label{eq:path-entry-expansion}
\end{equation}
The last factor accounts for the fractional delay in the waveform.
This expression corresponds to \cite[eq.~(5)]{huang2026} with normalized
Doppler $k_{\nu,\mathrm{Huang}}=-\nu NT_s$ under our sign convention.

For integer $d$, \eqref{eq:path-operator} reduces to
\begin{align}
 \mathbf G(d,\nu;g)
 &=\mathbf F_{\rm A}\mathbf D_{\nu,g}\bm\Omega_d
   \bm\Pi_{d\bmod N}\mathbf F_{\rm A}^{\mathsf H},
 \label{eq:integer-operator}\\
 [\bm\Omega_d]_{tt}
 &=e^{j2\pi c_1[(t-d)^2-\langle t-d\rangle_N^2]},
 \label{eq:integer-prefix-phase}
\end{align}
where $[\bm\Pi_l\mathbf s]_t=s[\langle t-l\rangle_N]$.
The full delay enters the prefix phase, and its remainder determines
the cyclic shift.

At the reference frame, the channel of terminal $k$ is given by
\begin{equation}
 \mathbf G_k=\sum_\ell g_{k\ell}
 \mathbf G(\tau_{k\ell}/T_s,\nu_k;N_{\rm CPP}),
 \qquad
 H_{kmrn}=[\mathbf a_k]_m[\mathbf G_k]_{rn},
 \label{eq:channel}
\end{equation}
where $g_{k\ell}$ includes the path gain after large-scale power
compensation. The tensor $H_k\in\mathbb C^{M\times N\times N}$ contains
all antenna, DAFT output, and DAFT input coordinates. Reshaping it as a matrix
in $\mathbb C^{M\times N^2}$ with one row per antenna yields a matrix of rank at most one under
the common spatial response assumption.

Pilots are formed from selected rows and permuted columns of a
$KN$-point discrete Fourier transform (DFT). With distinct row indices
$r_q$, a column permutation $c_{kn}$, and independent uniform phases
$\phi_{kn}$, set
\begin{equation}
 P_{qkn}=\sqrt{\frac{E_p}{N}}
 e^{-j2\pi r_qc_{kn}/(KN)}e^{j\phi_{kn}},
 \qquad \sum_{n=0}^{N-1}|P_{qkn}|^2=E_p.
 \label{eq:pilot}
\end{equation}
The common terminal Doppler introduces the phase increment
$\beta_k=\operatorname{wrap}(2\pi\nu_kT_{\rm frame})$ between consecutive
frames, where $\operatorname{wrap}$ maps an angle to $[-\pi,\pi)$.
The observation operator with these phase increments is
\begin{equation}
 A_{qkn}=P_{qkn}e^{jq\beta_k},\qquad q=0,\ldots,Q-1.
 \label{eq:effective-pilot}
\end{equation}

Define the effective channel $X_k=\alpha_kH_k$.
After prefix removal and DAFT demodulation, the received signal can be expressed as
\begin{equation}
 Y_{qmr}=\sum_{k=1}^K\sum_{n=0}^{N-1}A_{qkn}X_{kmrn}+W_{qmr},
 \qquad W_{qmr}\sim\CN(0,\sigma^2),
 \label{eq:observation-tensor}
\end{equation}
with conditionally independent noise entries.
Using the indexing
\begin{equation}
 [\mathbf Y]_{q,mN+r}=Y_{qmr},\quad
 [\mathbf A]_{q,(k-1)N+n}=A_{qkn},\quad
 [\mathbf X]_{(k-1)N+n,mN+r}=X_{kmrn},
 \label{eq:mmv-indexing}
\end{equation}
the multiple measurement vector (MMV) model can be written as
\begin{equation}
 \mathbf Y=\mathbf A\mathbf X+\mathbf W,\quad
 \mathbf Y\in\mathbb C^{Q\times MN},\quad
 \mathbf A\in\mathbb C^{Q\times KN},\quad
 \mathbf X\in\mathbb C^{KN\times MN}.
 \label{eq:mmv}
\end{equation}
All $MN$ measurement vectors share the terminal activity pattern.
The receiver takes the observation tensor, the original pilots $P$,
and the activity marginal $\rho$ as inputs. It estimates
$\widehat{\bm\beta}$ and $\widehat\sigma^2$ from these observations
and constructs $\widehat{\mathbf A}=\mathbf A(\widehat{\bm\beta})$.
The physical operator $\mathbf A=\mathbf A(\bm\beta)$ and the noise
variance $\sigma^2$ in \eqref{eq:mmv} describe data generation.
For simulations, the signal-to-noise ratio (SNR)
is defined using the clean superposition $\mathbf Y_0=\mathbf A\mathbf X$:
\begin{equation}
 \sigma^2=\frac{\norm{\mathbf Y_0}_F^2}{QMN}
 10^{-\mathrm{SNR}_{\rm dB}/10}.
 \label{eq:snr}
\end{equation}
The received signal power includes cross terms between terminals.

\subsection{Problem Formulation}

Let $D=MN^2$ and let $q_{\rm tr}$ be the distribution of normalized channels
generated from the physical model before applying activity masks.
We normalize channels and observations by a common scale $s$,
the root mean square of the training channels:
$(\mathbf Y,H_k,X_k,\mathbf W)$ are divided by $s$, and both
$\sigma^2$ and $\widehat\sigma^2$ by $s^2$. The pilot operators are
unchanged. The same symbols denote the normalized quantities below;
the reported channel estimates are rescaled by $s$.

For joint activity detection and channel estimation (JADCE), we use
the factorized prior
\begin{equation}
 q(\{dH_k\},\bm\alpha)=
 \prod_{k=1}^{K}q_{\rm tr}(dH_k)
 \rho^{\alpha_k}(1-\rho)^{1-\alpha_k}.
 \label{eq:working-prior}
\end{equation}
We approximate the joint prior by the product of its channel
and activity marginals. With the estimated phase and noise variance
held fixed, the approximate posterior distribution can be expressed as
\begin{equation}
 \begin{aligned}
 &q(d\mathbf X,\{dH_k\},\bm\alpha\mid
 \mathbf Y,\widehat{\mathbf A},\widehat\sigma^2)\\
 &\quad\propto
 e^{-\norm{\mathbf Y-\widehat{\mathbf A}\mathbf X}_F^2/\widehat\sigma^2}
 \prod_{k=1}^{K}
 \delta_{\alpha_kH_k}(dX_k)q_{\rm tr}(dH_k)
 \rho^{\alpha_k}(1-\rho)^{1-\alpha_k},
 \end{aligned}
 \label{eq:working-posterior}
\end{equation}
where $\delta_{\alpha_kH_k}$ enforces $X_k=\alpha_kH_k$.
Under this posterior, the minimum mean-squared error (MMSE) target
and the activity decision are given by
\begin{equation}
 \begin{aligned}
 X_k^{\rm MMSE}&=\E_q[X_k\mid\mathbf Y,\widehat{\mathbf A},\widehat\sigma^2],\\
 p_k^{\rm Bayes}&=q(\alpha_k=1\mid\mathbf Y,\widehat{\mathbf A},\widehat\sigma^2),\qquad
 \alpha_k^{\rm Bayes}=\mathbf1\{p_k^{\rm Bayes}\geq\tfrac12\}.
 \end{aligned}
 \label{eq:targets}
\end{equation}
The threshold corresponds to equal costs for a missed detection and a
false alarm. The receiver approximates these targets by iterating between
linear estimation and channel and activity denoising, returning
$\widehat X_k$, $p_k$, and $\widehat\alpha_k=\mathbf1\{p_k\geq1/2\}$.
